\documentclass[10pt,a4paper]{amsart}
\usepackage{amsmath,amssymb,amsthm,graphicx}
\usepackage{verbatim}
\newenvironment{reference}{\comment}{\endcomment}
%\newenvironment{reference}{}{}
\newenvironment{slogan}{\comment}{\endcomment}
\newenvironment{history}{\comment}{\endcomment}

% For commutative diagrams we use Xy-pic
\usepackage[all]{xy}

% We use 2cell for 2-commutative diagrams.
\xyoption{2cell}
% Theorem environments.
%
\theoremstyle{plain}
\newtheorem{theorem}[subsection]{Theorem}
\newtheorem{proposition}[subsection]{Proposition}
\newtheorem{lemma}[subsection]{Lemma}

\theoremstyle{definition}
\newtheorem{definition}[subsection]{Definition}
\newtheorem{example}[subsection]{Example}
\newtheorem{exercise}[subsection]{Exercise}
\newtheorem{situation}[subsection]{Situation}

\theoremstyle{remark}
\newtheorem{remark}[subsection]{Remark}
\newtheorem{remarks}[subsection]{Remarks}

\numberwithin{equation}{subsection}

% Macros
%
\def\lim{\mathop{\mathrm{lim}}\nolimits}
\def\colim{\mathop{\mathrm{colim}}\nolimits}
\def\Spec{\mathop{\mathrm{Spec}}}
\def\Hom{\mathop{\mathrm{Hom}}\nolimits}
\def\Ext{\mathop{\mathrm{Ext}}\nolimits}
\def\SheafHom{\mathop{\mathcal{H}\!\mathit{om}}\nolimits}
\def\SheafExt{\mathop{\mathcal{E}\!\mathit{xt}}\nolimits}
\def\Sch{\mathit{Sch}}
\def\Mor{\mathop{\mathrm{Mor}}\nolimits}
\def\Ob{\mathop{\mathrm{Ob}}\nolimits}
\def\Sh{\mathop{\mathit{Sh}}\nolimits}
\def\NL{\mathop{N\!L}\nolimits}
\def\CH{\mathop{\mathrm{CH}}\nolimits}
\def\proetale{{pro\text{-}\acute{e}tale}}
\def\etale{{\acute{e}tale}}
\def\QCoh{\mathit{QCoh}}
\def\Ker{\mathop{\mathrm{Ker}}}
\def\Im{\mathop{\mathrm{Im}}}
\def\Coker{\mathop{\mathrm{Coker}}}
\def\Coim{\mathop{\mathrm{Coim}}}

% Boxtimes
%
\DeclareMathSymbol{\boxtimes}{\mathbin}{AMSa}{"02}

%
% Macros for moduli stacks/spaces
%
\def\QCohstack{\mathcal{QC}\!\mathit{oh}}
\def\Cohstack{\mathcal{C}\!\mathit{oh}}
\def\Spacesstack{\mathcal{S}\!\mathit{paces}}
\def\Quotfunctor{\mathrm{Quot}}
\def\Hilbfunctor{\mathrm{Hilb}}
\def\Curvesstack{\mathcal{C}\!\mathit{urves}}
\def\Polarizedstack{\mathcal{P}\!\mathit{olarized}}
\def\Complexesstack{\mathcal{C}\!\mathit{omplexes}}
% \Pic is the operator that assigns to X its picard group, usage \Pic(X)
% \Picardstack_{X/B} denotes the Picard stack of X over B
% \Picardfunctor_{X/B} denotes the Picard functor of X over B
\def\Pic{\mathop{\mathrm{Pic}}\nolimits}
\def\Picardstack{\mathcal{P}\!\mathit{ic}}
\def\Picardfunctor{\mathrm{Pic}}
\def\Deformationcategory{\mathcal{D}\!\mathit{ef}}
\usepackage{cleveref}
\crefname{originalreference}{original reference}{original references}
\Crefname{originalreference}{Original reference}{Original references}
\makeatletter
\newlabel{spaces-topologies-lemma-fpqc}{{03MR}{source}{}{sourcepointer-1}{}}
\newlabel{spaces-topologies-lemma-fpqc@cref}{{[originalreference][0][]03MR}{[1][0][]source}{}{}{}}
\makeatother
\crefname{originallemma}{original lemma}{original lemmas}
\Crefname{originallemma}{Original lemma}{Original lemmas}
\makeatletter
\newlabel{lemma-descending-property-isomorphism}{{041Y}{source}{}{sourcepointer-2}{}}
\newlabel{lemma-descending-property-isomorphism@cref}{{[originallemma][0][]041Y}{[1][0][]source}{}{}{}}
\makeatother
\crefname{originalreference}{original reference}{original references}
\Crefname{originalreference}{Original reference}{Original references}
\makeatletter
\newlabel{spaces-topologies-lemma-refine-fpqc-schemes}{{0419}{source}{}{sourcepointer-3}{}}
\newlabel{spaces-topologies-lemma-refine-fpqc-schemes@cref}{{[originalreference][0][]0419}{[1][0][]source}{}{}{}}
\makeatother
\crefname{originalreference}{original reference}{original references}
\Crefname{originalreference}{Original reference}{Original references}
\makeatletter
\newlabel{spaces-properties-proposition-sheaf-fpqc}{{0APL}{source}{}{sourcepointer-4}{}}
\newlabel{spaces-properties-proposition-sheaf-fpqc@cref}{{[originalreference][0][]0APL}{[1][0][]source}{}{}{}}
\makeatother
\title{Being a monomorphism is fpqc local on the base (two proofs: diagonal descent and reduction to schemes)}
\author{The Stacks Project Authors}
\date{stacks-project:spaces-descent.tex:local-snapshot-20261012; GNU FDL 1.2 or later}
\begin{document}
\maketitle
\noindent\textit{Editorial note.} All delivered mathematical text is the author's own, copied verbatim from the pinned Stacks Project chapter file spaces-descent.tex; no statement or proof step was rewritten, reordered or removed.\par\medskip
\noindent\textit{Editorial note.} The counted span contains the author's two labelled arguments 'First proof.' and 'Second proof.', kept verbatim; both are complete derivations and neither is a pointer.\par\medskip
\noindent\textit{Editorial note.} The counted span cites two Topologies on Spaces lemmas, one same-chapter lemma and one Properties of Spaces proposition, recorded in assembly.external\_references with canonical tags as navigation pointers.\par\medskip
\noindent\textit{Allowed standard prerequisites (editorial).}\begin{itemize}
\item Monomorphism iff diagonal is an isomorphism. Standard, used by the author inside the proof. Supplies the reformulation driving the first proof.
\item fpqc coverings are stable under base change. Stacks Project, Topologies on Spaces, Lemma spaces-topologies-lemma-fpqc. Supplies the covering of X \textbackslash{}times\_Y X used in the descent step.
\item fpqc descent of the isomorphism property. Stacks Project, Descent on Spaces, Lemma lemma-descending-property-isomorphism. Supplies that \textbackslash{}Delta is an isomorphism when its base changes are.
\item Refinement of fpqc coverings by morphisms of schemes. Stacks Project, Topologies on Spaces, Lemma spaces-topologies-lemma-refine-fpqc-schemes. Supplies the reduction of the covering to schemes in the second proof.
\item Sheaves form an fpqc sheaf / local equality detection. Stacks Project, Properties of Spaces, Proposition spaces-properties-proposition-sheaf-fpqc. Supplies a = b from agreement over the fpqc covering.
\end{itemize}\medskip
\ifdefined\hypertarget\hypertarget{sourcepointer-1}{}\fi
\noindent\textit{Source reference.} Original reference 03MR: sources/stacks-chow/tags.txt, source lines 4759--4759. Cross-chapter standard prerequisite cited by the author inside the counted span; it is used as a pointer and its own proof is not part of this item.\par\medskip
\ifdefined\hypertarget\hypertarget{sourcepointer-2}{}\fi
\noindent\textit{Source reference.} Original lemma 041Y: sources/stacks-chow/tags.txt, source lines 5309--5309. Cross-chapter standard prerequisite cited by the author inside the counted span; it is used as a pointer and its own proof is not part of this item.\par\medskip
\ifdefined\hypertarget\hypertarget{sourcepointer-3}{}\fi
\noindent\textit{Source reference.} Original reference 0419: sources/stacks-chow/tags.txt, source lines 5284--5284. Cross-chapter standard prerequisite cited by the author inside the counted span; it is used as a pointer and its own proof is not part of this item.\par\medskip
\ifdefined\hypertarget\hypertarget{sourcepointer-4}{}\fi
\noindent\textit{Source reference.} Original reference 0APL: sources/stacks-chow/tags.txt, source lines 13479--13479. Cross-chapter standard prerequisite cited by the author inside the counted span; it is used as a pointer and its own proof is not part of this item.\par\medskip
\begin{lemma}
\label{lemma-descending-property-monomorphism}
The property $\mathcal{P}(f) =$``$f$ is a monomorphism''
is fpqc local on the base.
\end{lemma}

\begin{proof}
Let $f : X \to Y$ be a morphism of algebraic spaces.
Let $\{Y_i \to Y\}$ be an fpqc covering, and assume
each of the base changes $f_i : X_i \to Y_i$ of $f$ is
a monomorphism. We have to show that $f$ is a monomorphism.

\medskip\noindent
First proof. Note that $f$ is a monomorphism if and only if
$\Delta : X \to X \times_Y X$ is an isomorphism. By applying this to
$f_i$ we see that each of the morphisms
$$
\Delta_i :
X_i
\longrightarrow
X_i \times_{Y_i} X_i = Y_i \times_Y (X \times_Y X)
$$
is an isomorphism. The base change of an fpqc covering is an fpqc covering, see
Topologies on Spaces, Lemma \ref{spaces-topologies-lemma-fpqc}
hence $\{Y_i \times_Y (X \times_Y X) \to X \times_Y X\}$
is an fpqc covering of algebraic spaces. Moreover, each
$\Delta_i$ is the base change of the morphism
$\Delta : X \to X \times_Y X$. Hence it follows from
Lemma \ref{lemma-descending-property-isomorphism}
that $\Delta$ is an isomorphism, i.e., $f$ is a monomorphism.

\medskip\noindent
Second proof.
Let $V$ be a scheme, and let $V \to Y$ be a surjective \'etale morphism.
If we can show that $V \times_Y X \to V$ is a monomorphism, then it
follows that $X \to Y$ is a monomorphism. Namely, given any
cartesian diagram of sheaves
$$
\vcenter{
\xymatrix{
\mathcal{F} \ar[r]_a \ar[d]_b & \mathcal{G} \ar[d]^c \\
\mathcal{H} \ar[r]^d & \mathcal{I}
}
}
\quad
\quad
\mathcal{F} = \mathcal{H} \times_\mathcal{I} \mathcal{G}
$$
if $c$ is a surjection of sheaves, and $a$ is injective, then also
$d$ is injective. This reduces the problem to the case where $Y$ is
a scheme. Moreover, in this case we may assume that the algebraic spaces
$Y_i$ are schemes also, since we can always refine the covering to place
ourselves in this situation, see
Topologies on Spaces, Lemma \ref{spaces-topologies-lemma-refine-fpqc-schemes}.

\medskip\noindent
Assume $\{Y_i \to Y\}$ is an fpqc covering of schemes.
Let $a, b : T \to X$ be two morphisms
such that $f \circ a = f \circ b$. We have to show that $a = b$.
Since $f_i$ is a monomorphism we see that $a_i = b_i$, where
$a_i, b_i : Y_i \times_Y T \to X_i$ are
the base changes. In particular the compositions
$Y_i \times_Y T \to T \to X$ are equal.
Since $\{Y_i \times_Y T \to T\}$ is an fpqc covering we
deduce that $a = b$ from Properties of Spaces, Proposition
\ref{spaces-properties-proposition-sheaf-fpqc}.
\end{proof}

\end{document}
